%!TEx root=../a03.tex
\section{Q4}
\begin{enumerate}[label = (\alph*)]
    \item \begin{solution}
        This statement is true. 

        Suppose $x$ be arbitrary in $\mathcal{F}$. By the definition of set, we have 
        \begin{equation*}
            42 \mid x + 18, x \in \Z. 
        \end{equation*} 

        By the definition of divisibility, there exists $k \in \Z$ such that 
        \begin{equation*}
            42k = x + 18. 
        \end{equation*}

        Rearrange, we get 
        \begin{equation*}
            x = 42k - 18. 
        \end{equation*}

        Factor 3 out, we get 
        \begin{equation*}
            x = 3\paren{14k - 6}.
        \end{equation*}

        Since $k \in \Z$, 
        \begin{equation*}
            x = 3\paren{14k - 6}, 14k - 6 \in \Z. 
        \end{equation*}

        This yields $x \in \mathcal{K}$. 

        Moreover, substitute $x = 42k - 18$ into $x + 18$, we get 
        \begin{equation*}
            42k - 18 + 18 = 42k. 
        \end{equation*}

        This is equal to 
        \begin{equation*}
            7 \cdot \paren{6k}. 
        \end{equation*}

        Since $k \in \Z$, $6k \in \Z$. 

        Therefore, there exists $6k \in \Z$ such that 
        \begin{equation*}
            x + 18 = 7 \cdot 6k. 
        \end{equation*}

        By the definition of divisibility, we have 
        \begin{equation*}
            7 \mid x + 18. 
        \end{equation*}

        We obtain $x \in \mathcal{P}$. 

        Therefore, we proved 
        \begin{equation*}
            \forall x \in \Z, \bracks{\paren{x \in \mathcal{F}} \implies \paren{x \in \paren{\mathcal{K} \cap \mathcal{P}}}}
        \end{equation*}
    \end{solution}

    \newpage
    \item \begin{solution}
        This statement is false. 

        We disprove this statement by giveing a counterexample. Let $x = 45 \in \Z$. 
        
        Since 
        \begin{equation*}
            45 + 18 = 63
        \end{equation*}
        and 
        \begin{equation*}
            42 \nmid 63. 
        \end{equation*}
        We have 
        \begin{equation*}
            45 \notin \mathcal{F}
        \end{equation*}

        However, $3 \mid 45$. 
        So 
        \begin{equation*}
            45 \in \mathcal{K}. 
        \end{equation*}

        Moreover, 
        \begin{equation*}
            45 + 18 = 63. 
        \end{equation*}

        and 
        \begin{equation*}
            7 \mid 63. 
        \end{equation*}

        So 
        \begin{equation*}
            45 \in \mathcal{P}. 
        \end{equation*}

        Thus, we proved that 
        \begin{equation*}
            \exists x \in \Z, \bracks{\paren{x \notin \mathcal{F}} \wedge \paren{x \in \paren{\mathcal{K} \cap \mathcal{P}}}}
        \end{equation*}
        , which is a negation of the original statement. 
        So, we disprove the original statement. 
    \end{solution}
\end{enumerate}